\section{精确对称类型与渐近行为}\label{sec:types}

\begin{theorem}[精确稳定子分类]\label{thm:stabilizers}
对三行、$2n$ 列棋盘，各精确稳定子的铺法数和轨道数见表 \ref{tab:stabilizers}。
该表对所有 $n\ge0$ 成立。
\end{theorem}
\begin{table}[htbp]
\centering
\caption{三行铺法的全部精确稳定子类型。}\label{tab:stabilizers}
\begin{tabular}{cccc}
\toprule
精确稳定子 & 铺法数 & 轨道大小 & 轨道数\\
\midrule
$G$ & $1$ & $1$ & $1$\\
$\{e,h\}$ & $0$ & $2$ & $0$\\
$\{e,v\}$ & $v_n-1$ & $2$ & $(v_n-1)/2$\\
$\{e,\rho\}$ & $r_n-1$ & $2$ & $(r_n-1)/2$\\
$\{e\}$ & $d_n-v_n-r_n+1$ & $4$ & $(d_n-v_n-r_n+1)/4$\\
\bottomrule
\end{tabular}
\end{table}
\begin{proof}
引理 \ref{lem:h-odd} 表明唯一的上下反射固定铺法是三行各自水平铺满的铺法，
它同时被 $v,\rho$ 固定。
若铺法同时由任意两个不同的非恒等群元素固定，则也由它们的积固定，
因而由整个 $G$ 固定。故三种固定集的两两交集和三重交集都是上述单点。
于是只具 $v$ 对称的铺法有 $v_n-1$ 个，只具 $\rho$ 对称的有 $r_n-1$ 个，
只具 $h$ 对称的有零个；其余 $d_n-v_n-r_n+1$ 个的稳定子平凡。
群 $G$ 交换，同一轨道中的稳定子相同。各类除以轨道大小即得表中结果。
四元群的子群只有表中五种，故分类完整。
\end{proof}

例如 $n=8$ 时 $d_8=29681,v_8=265,r_8=153$。
共有 $7525$ 个轨道，其中 $7316$ 个稳定子平凡，
$132$ 个精确稳定子为 $\{e,v\}$，$76$ 个为 $\{e,\rho\}$，$1$ 个为 $G$。

\begin{theorem}[对称铺法的渐近比例]\label{thm:asym}
令 $\lambda=2+\sqrt3$。均匀选取一个有固定棋盘方向的普通铺法时，
其具有非平凡对称的概率 $p_n$ 满足
\begin{align}
 p_{2j}&\sim(\lambda-1)\lambda^{-j},&
 p_{2j+1}&\sim\frac{\lambda+1}{\lambda}\lambda^{-j}
 \quad(j\to\infty).\label{eq:prob-asym}
\end{align}
特别地，$p_n=O(\lambda^{-n/2})\to0$，并且
\begin{equation}\label{eq:orbit-asym}
 a(n)=\frac{d_n}{4}+O(\lambda^{n/2})
 \sim\frac{3+\sqrt3}{24}\lambda^n.
\end{equation}
\end{theorem}
\begin{proof}
由 \eqref{eq:d-rec} 的特征根和初值，
\begin{equation}\label{eq:d-closed}
 d_n=c_+\lambda^n+c_-\lambda^{-n},\qquad
 c_+=\frac{3+\sqrt3}{6},\quad c_-=\frac{3-\sqrt3}{6}.
\end{equation}
表 \ref{tab:stabilizers} 给出
$p_n=(v_n+r_n-1)/d_n$。
而 \eqref{eq:parity-fixed} 给出
\[
 v_{2j}+r_{2j}=d_{j+1}-d_j,\qquad
 v_{2j+1}+r_{2j+1}=d_{j+1}+d_j.
\]
将 \eqref{eq:d-closed} 代入，分子主项分别为
$c_+(\lambda-1)\lambda^j$ 与 $c_+(\lambda+1)\lambda^j$，
分母主项分别为 $c_+\lambda^{2j}$ 与 $c_+\lambda^{2j+1}$，
故得 \eqref{eq:prob-asym}。
最后由 $a(n)=(d_n+1+v_n+r_n)/4$ 得到 \eqref{eq:orbit-asym}。
\end{proof}

如果均匀选取的是轨道，则非平凡对称轨道数为
$1+(v_n-1)/2+(r_n-1)/2=(v_n+r_n)/2$，
其比例为 $(v_n+r_n)/(2a(n))$，同样趋于零，但这是不同的概率空间。
